% !TeX root = main.tex
\chapter[Federated Training Algorithms]{FL Algorithms}
\label{appendix:algorithms}

This appendix gives the federated learning methods used in the dissertation in a unified pseudocode style. It is intended for comparison and uses the notation collected in Appendix~\ref{chap:notation}. The pseudocode follows the implementation behaviour, with logging, centralised evaluation, and artefact collection omitted.


\section{Full-model synchronous training}

\begin{algorithm}[H]
\caption{\textcolor{fedavggreen}{\texttt{FedAvg}}}
\label{alg:fedavg}
\begin{algorithmic}[1]
\Require total clients $\mathcal{K}$, synchronous rounds $T$, local learning rate $\eta_\ell$, local epochs $E$
\Ensure final global model $w^{(T)}$
\Statex \textbf{Server initialisation:}
\State initialise global model $w^{(0)}$
\For{$t = 0$ to $T-1$}
    \State sample client set $\mathcal{C}_t \subseteq \mathcal{K}$
    \ForAll{$k \in \mathcal{C}_t$ \textbf{in parallel}}
        \State send $w^{(t)}$ to client $k$
    \EndFor

    \Statex
    \Statex \hspace{1em}\textbf{Client $k \in \mathcal{C}_t$:}
    \State initialise local model $w_k \gets w^{(t)}$
    \For{local epoch $e = 1$ to $E$}
        \State update $w_k$ with local SGD on client $k$'s data
    \EndFor
    \State return $\tilde{w}_k^{(t+1)} \gets w_k$

    \Statex
    \Statex \textbf{Server aggregation:}
    \State compute example-weighted mean over $\mathcal{C}_t$
    \[
        \bar{\Delta}^{(t)} \gets
        \frac{\sum_{k \in \mathcal{C}_t} n_k\left(\tilde{w}_k^{(t+1)} - w^{(t)}\right)}
             {\sum_{h \in \mathcal{C}_t} n_h}
    \]
    \State \textcolor{fedavggreen}{\textbf{FedAvg:} \hspace{0.4em} $w^{(t+1)} \gets w^{(t)} + \bar{\Delta}^{(t)}$}
\EndFor
\end{algorithmic}
\end{algorithm}

\section{Selective partial-training methods}

The three implemented submodel FL methods share the same overall structure: the server defines a client-visible trainable subset, clients train only that subset locally, and the server selectively aggregates only the parameters that each client actually updated. The methods differ in how the trainable subset is constructed and in how returned updates are combined.

\begin{algorithm}[H]
\caption{\textcolor{heteroflblue}{\texttt{HeteroFL}}, \textcolor{fedrolexpurple}{\texttt{FedRolex}}, and \textcolor{fiarsered}{\texttt{FIARSE}}}
\label{alg:partial_training_family}
\begin{algorithmic}[1]
\Require total clients $\mathcal{K}$, local example counts $\{n_k\}$, rounds $T$, local epochs $E$, model-rate lookup $\mathcal{A}$, extraction policy $\Psi$
\Ensure final global model $w^{(T)}$
\Statex \textbf{Server initialisation:}
\State initialise full global model $w^{(0)}$
\For{$t = 0$ to $T-1$}
    \State sample client set $\mathcal{C}_t \subseteq \mathcal{K}$
    \State look up fixed model rates $\{r_k : k \in \mathcal{C}_t\} \gets \mathcal{A}(\mathcal{C}_t)$
    \ForAll{$k \in \mathcal{C}_t$ \textbf{in parallel}}
        \State compute trainable subset $I_k^{(t)} \gets \Psi(w^{(t)}, r_k, t)$ \Comment{\textcolor{heteroflblue}{static nested} / \textcolor{fedrolexpurple}{rolling window} / \textcolor{fiarsered}{sparse mask}}
        \State construct client-visible model $w_{k,0}^{(t)}$ from $w^{(t)}$ and $I_k^{(t)}$
        \State send $w_{k,0}^{(t)}$ to client $k$
    \EndFor

    \Statex
    \Statex \hspace{1em}\textbf{Client $k \in \mathcal{C}_t$:}
    \State initialise local model $w_k \gets w_{k,0}^{(t)}$
    \For{local epoch $e = 1$ to $E$}
        \State update $w_k$ with local SGD on client $k$'s data
    \EndFor
    \State return $\tilde{w}_k^{(t+1)} \gets w_k$

    \Statex
    \Statex \textbf{Server selective aggregation:}
    \For{each global parameter index $j$}
        \State $\mathcal{U}_j^{(t)} \gets \{k \in \mathcal{C}_t : j \in I_k^{(t)}\}$
        \If{$\mathcal{U}_j^{(t)} \neq \emptyset$}
            \State $[w^{(t+1)}]_j \gets
            \dfrac{\sum_{k \in \mathcal{U}_j^{(t)}} n_k [\tilde{w}_k^{(t+1)}]_j}
                  {\sum_{h \in \mathcal{U}_j^{(t)}} n_h}$ \Comment{selective \texttt{FedAvg} over contributors}
        \Else
            \State $[w^{(t+1)}]_j \gets [w^{(t)}]_j$
        \EndIf
    \EndFor
\EndFor
\end{algorithmic}
\end{algorithm}

\subsection{Method instantiations}

Within the shared partial-training skeleton, the methods differ primarily in how they define the trainable subset. \textcolor{heteroflblue}{\texttt{HeteroFL}} uses a static nested extractor with example-weighted selective averaging, \textcolor{fedrolexpurple}{\texttt{FedRolex}} uses a rolling extractor and defaults to equal client weighting in the reference paper, and \textcolor{fiarsered}{\texttt{FIARSE}} uses sparse importance masks with nonzero-delta aggregation.

\paragraph{\textcolor{heteroflblue}{\texttt{HeteroFL}}.}
The extraction policy $\Psi$ uses a \textbf{static nested} selection rule: for each hidden layer, the client receives a prefix subnetwork at rate $r_k$, and the server uses example-weighted selective averaging over the clients that contributed each parameter.

\paragraph{\textcolor{fedrolexpurple}{\texttt{FedRolex}}.}
\texttt{FedRolex} replaces static extraction with a \textbf{rolling window} over eligible channels, so the selected submodel advances with round index $t$. In the implemented variant, selective averaging is \textbf{unweighted}, matching the reference paper's default equal-client weighting.

\paragraph{\textcolor{fiarsered}{\texttt{FIARSE}}}.
\texttt{FIARSE} replaces dense slicing with an \textbf{importance-aware sparse mask}. For each maskable parameter \(j\), the client computes the salience score
\[
    s_j^{(t)} \gets |[w^{(t)}]_j|,
\]
then keeps parameters above the model-rate threshold. Two threshold modes are implemented:
\begin{itemize}[topsep=4pt, itemsep=4pt]
    \item \textbf{Layerwise:} compute a separate top \(r_k\) fraction threshold within each maskable parameter tensor.
    \item \textbf{Global:} compute one top \(r_k\) fraction threshold across all maskable parameter tensors.
\end{itemize}
During local training, the binary mask is applied with the straight-through estimator used by the FIARSE reference mechanism. At aggregation, the server averages only nonzero returned deltas for each parameter and applies the configured server learning rate.

\section{Buffered asynchronous training}

\begin{algorithm}[H]
\caption{\textcolor{fedbuffteal}{\texttt{FedBuff}} and \textcolor{fedstalebrown}{\texttt{FedStaleWeight}}}
\label{alg:fedbuff}
\begin{algorithmic}[1]
\Require available clients $\mathcal{K}$, asynchronous server steps $S$, train concurrency $M$, buffer size $K$, local learning rate $\eta_\ell$, server learning rate $\eta_s$
\Ensure final global model $w^{[S]}$
\Statex \textbf{Server initialisation:}
\State initialise global model $w^{[0]}$
\State initialise buffered update set $\mathcal{B} \gets \emptyset$
\State initialise in-flight request set $\mathcal{Q} \gets \emptyset$
\State dispatch up to $M$ train requests tagged with model version $0$
\State set server step $s \gets 0$
\While{$s < S$}
    \State pull any replies that have arrived from requests in $\mathcal{Q}$
    \ForAll{valid arriving replies $i$}
        \State let $v_i$ be the model version originally sent to that client
        \State let $\tilde{w}_i$ be the returned local model
        \State compute local delta $\Delta_i \gets w^{[v_i]} - \tilde{w}_i$
        \State compute staleness $\tau_i \gets s - v_i$
        \State append $(\Delta_i,\tau_i)$ to $\mathcal{B}$
        \State dispatch another train request if an idle client is available and $|\mathcal{Q}| < M$
    \EndFor
    \If{$|\mathcal{B}| = K$}
        \State choose buffer weights $\{\alpha_i : i \in \mathcal{B}\}$ \Comment{\textcolor{fedbuffteal}{uniform/decay} or \textcolor{fedstalebrown}{fair}}
        \State compute buffered update
        \[
            \bar{\Delta}^{[s]} \gets \sum_{i \in \mathcal{B}} \alpha_i \Delta_i
        \]
        \State apply server update $w^{[s+1]} \gets w^{[s]} - \eta_s \bar{\Delta}^{[s]}$
        \State clear $\mathcal{B}$
        \State increment server step $s \gets s + 1$
        \State retain only the historical model versions still referenced by in-flight requests
    \EndIf
\EndWhile
\end{algorithmic}
\end{algorithm}

\subsection{Implemented weighting modes for buffered asynchronous aggregation}

The dissertation implementation supports three weighting modes for buffered async updates:
\begin{align*}
    \textcolor{fedbuffteal}{\texttt{none}}: \qquad
    &g(\tau_i) = 1,
    &\alpha_i = \frac{g(\tau_i)}{\sum_{u \in \mathcal{B}} g(\tau_u)} = \frac{1}{K}, \\
    \textcolor{fedbuffteal}{\texttt{polynomial}}: \qquad
    &g(\tau_i) = (1 + \tau_i)^{-a},
    &\alpha_i = \frac{g(\tau_i)}{\sum_{u \in \mathcal{B}} g(\tau_u)}, \\
    \textcolor{fedstalebrown}{\texttt{fair}}: \qquad
    &g(\tau_i) = \tau_i K + 1,
    &\alpha_i = \frac{g(\tau_i)}{\sum_{u \in \mathcal{B}} g(\tau_u)}.
\end{align*}

The final mode is the fairness weighting used by the implemented \texttt{FedStaleWeight} variant. It increases the effective contribution of staler buffered updates and is analysed as a separate fairness-oriented asynchronous method built on the same buffered execution core.

\section{Coverage-aware asynchronous partial training}

\begin{algorithm}[H]
% \footnotesize
\caption{\texttt{FedCover}}
\label{alg:fedcover}
\begin{algorithmic}[1]
\Require available clients $\mathcal{K}$, asynchronous server steps $S$, train concurrency $M$, buffer size $K$, model-rate lookup $\mathcal{A}$, nested extractor $\Psi$, server learning rate $\eta_s$, coverage power $\gamma$, maximum gain $\kappa_{\max}$
\Ensure final global model $w^{[S]}$
\State initialise full global model $w^{[0]}$
\For{$s = 0,\ldots,S-1$}
    \State keep at most $M$ asynchronous client trainings in flight using fixed rates $r_k \gets \mathcal{A}(k)$ and slices $I_k \gets \Psi(w^{[s]}, r_k, s)$
    \State collect the next buffer $\mathcal{B}_s$ of $K$ valid replies
    \ForAll{replies $i \in \mathcal{B}_s$}
        \State recover the sent version $v_i$, trained index set $I_i$, and returned submodel $\tilde{w}_i$
        \State compute slice delta $\Delta_i \gets \textsc{Slice}(w^{[v_i]}, I_i)-\tilde{w}_i$ and staleness $\tau_i \gets s-v_i$
        \State set staleness score $q_i \gets g(\tau_i)$
    \EndFor
    \State normalise buffer weights $\alpha_i \gets q_i / \sum_{u \in \mathcal{B}_s} q_u$
    \State set $w^{[s+1]} \gets w^{[s]}$
    \ForAll{parameter indices $j \in \bigcup_{i\in\mathcal{B}_s} I_i$}
        \State compute partial delta and coverage mass
        \[
            [\bar{\Delta}^{[s]}]_j =
            \sum_{i\in\mathcal{B}_s}\alpha_i\mathbf{1}[j\in I_i][\Delta_i]_j,
            \qquad
            c_j^{[s]} =
            \sum_{i\in\mathcal{B}_s}\alpha_i\mathbf{1}[j\in I_i]
        \]
        \State set $\kappa_j^{[s]} \gets \min\!\left(\kappa_{\max}, (c_j^{[s]})^{-\gamma}\right)$
        \State update $[w^{[s+1]}]_j \gets [w^{[s]}]_j - \eta_s \kappa_j^{[s]} [\bar{\Delta}^{[s]}]_j$
    \EndFor
\EndFor
\end{algorithmic}
\end{algorithm}

\texttt{FedCover} differs from uncorrected \texttt{Async-HeteroFL} only in the coverage-gain step. With $\gamma=0$ or $\kappa_{\max}=1$, the coverage gain is one and the method reduces to buffered asynchronous partial training with the same fixed model rates. With full-model clients, every trained parameter has $c_j=1$, so the coverage correction again disappears.
